\subsection{Verlinde Reduction and Gauss Sum}

We now apply the Reshetikhin--Turaev (RT) construction to the framed braid \(\beta\) obtained in Section~\ref{sec:braid}. Recall from Section~\ref{sec:mtc} that \(\mathcal{C}_p = \mathrm{SU}(2)_{p-2}\) is a modular tensor category with simple objects labeled by \(j = 0, 1/2, \dots, (p-2)/2\), quantum dimensions
\[
d_j(p) = \frac{\sin\bigl(\pi(2j+1)/p\bigr)}{\sin(\pi/p)},
\]
total quantum dimension
\[
D_p = \sqrt{\sum_j d_j(p)^2} = \sqrt{\frac{p}{2\sin^2(\pi/p)}},
\]
S-matrix entries
\[
S_{jj'} = \sqrt{\frac{2}{p}} \sin\biggl( \frac{\pi(2j+1)(2j'+1)}{p} \biggr),
\]
and ribbon element (framing phase)
\[
\theta_j(p) = \exp\biggl( \frac{\pi i}{p} \bigl( j(j+1) - \tfrac14 \bigr) \biggr).
\]
The T-matrix is the diagonal matrix with entries \(\theta_j(p)\). The Verlinde formula expresses the fusion coefficients, but for the RT invariant of a closed oriented 3-manifold \(M\) presented as the closure of a framed braid (or equivalently as a Seifert fibered space), the invariant takes the explicit form (see Turaev \cite{TuraevBook} and Kirby--Melvin \cite{KirbyMelvin} for the general Seifert case)
\[
\mathrm{RT}_{\mathcal{C}_p}(M) = \frac{1}{D_p^{2g-2 + r}} \sum_{j_1,\dots,j_r} \Biggl( \prod_{\ell=1}^r d_{j_\ell} \Biggr) \, S_{0j_1} \cdots S_{0j_r} \, \prod_{k=1}^s \theta_{j_k}^{m_k} \, \exp\bigl(2\pi i \, e \, h\bigr),
\]
where \(g\) is the genus of the base orbifold, \(r\) the number of exceptional fibers, \(m_k\) the twisting exponents, \(e\) the Euler number of the Seifert fibration, and \(h\) is a linear combination of the quadratic Casimir values \(j_k(j_k+1)\). (The precise exponents depend on the Seifert invariants \((\alpha_i,\beta_i)\) and the framing of the braid; full normalization conventions follow the ``corrected'' S-matrix given above.)

In our setting the closure of \(\beta\) (on \(n\) strands) is claimed to be a Seifert fibered space over a 4-punctured sphere (or equivalently a mapping torus of a periodic orbit). Assuming the braid encoding of Section~\ref{sec:braid} yields explicit Seifert invariants such that the state sum collapses after repeated application of the Verlinde fusion rules and orthogonality of the S-matrix, the invariant formally reduces to a single matrix element in the \(j'=0\) column:
\[
\mathrm{RT}_{\mathcal{C}_p}(\beta) = \frac{1}{D_p^{n-1}} \, \Biggl( \prod_{k=1}^{n} \theta_{j_k}^{e_k} \Biggr) \, G_p,
\]
where the prefactor \(1/D_p^{n-1}\) arises from the \(n-1\) independent colorings after fixing the trivial representation on one component (standard for braid closures), the product of framing phases \(\theta_{j_k}^{e_k}\) collects the blackboard framing correction plus the writhe of the braid word, and the reduced sum is the Gauss sum
\[
G_p := \sum_{j} d_j(p) \, \theta_j(p) \, S_{j0}.
\]
(The precise exponents \(e_k\) are determined by the canonical framing prescription of Section~\ref{sec:braid}.)

A direct symbolic evaluation of \(G_p\) using the explicit formulas above yields, for the smallest primes,
\begin{align*}
G_3 &= \frac{\sqrt{3}}{2} \exp(-i\pi/12) \cdot (\text{explicit cyclotomic factor}), \\
G_5 &= \text{lengthy linear combination of } \sqrt{5\pm\sqrt{5}} \text{ and roots of unity (not equal to } e^{i\pi/4}5^{-1/2}\text{)}.
\end{align*}
In general \(G_p \neq \tau_p\) with \(\tau_p = e^{\pi i/4} p^{-1/2}\). The classical Landsberg--Schaar reciprocity applies only to unweighted quadratic Gauss sums \(\sum_{k=0}^{p-1} \exp(2\pi i a k^2/p)\), whereas \(G_p\) is a weighted sum involving quantum dimensions \(d_j(p)\) and the trigonometric S-matrix entries; no identity in the literature converts the weighted sum into the required \(\tau_p\).

Thus the formal reduction above does \emph{not} yield \(\mathrm{RT}_{\mathcal{C}_p}(\beta) = p^{-n}\) after cancellation of phases and normalizations.

\begin{remark}[Gaps to be bridged or left unbridgeable]
The step cannot be completed rigorously for the following independent reasons:

1. \textbf{Missing Seifert data.} The braid word \(\beta \in B_n\) produced by the tree shift (Phase~2) has not been shown to close to a Seifert manifold whose invariants \((\alpha_i,\beta_i,e)\) make the Verlinde sum collapse exactly to the single Gauss sum \(G_p\) (with the stated prefactors). Without the explicit Seifert invariants, the reduction formula itself remains conjectural.

2. \textbf{Falsity of the Gauss-sum identity.} Even granting the reduction, the key claim \(G_p = \tau_p\) (up to the overall \(p^{-n}\) factor) is false, as verified by exact cyclotomic arithmetic (SymPy implementation over \(\mathbb{Q}(\zeta_p)\), confirming mismatch of both magnitude and argument for all \(p\leq 17\)). Landsberg--Schaar reciprocity cannot be applied directly because the sum is not a classical quadratic Gauss sum.

3. \textbf{Framing and normalization bookkeeping.} While the blackboard framing correction and ribbon phases can be written explicitly once \(\beta\) is given, their cancellation to a pure power \(p^{-n}\) relies on the erroneous Gauss-sum evaluation.

These gaps cannot be bridged within the current construction without either (a) a fundamentally different braid-to-manifold identification that alters the Verlinde matrix element, or (b) a modified definition of the RT invariant (outside standard quantum topology). The local identity \(\mathrm{RT}_{\mathcal{C}_p}(\beta)=p^{-n}\) therefore remains an open conjecture rather than a theorem. All subsequent global adelic constructions that rely on this local weight are correspondingly conditional.

A complete, reproducible verification suite (Python/SageMath script computing the full Verlinde sum for any candidate braid word) is provided in Appendix~C; it confirms the mismatch for all tested periodic orbits.
\end{remark}

The remainder of the proof (Phase~3.3 onward) therefore cannot proceed as written. We retain the above formulas as a precise record of the attempted reduction, which may serve as a starting point for future corrections of the braid encoding or the MTC data.